\begin{abstract}
This dissertation develops statistical methods and computational tools for
conditional quantiles and fixed-content intervals in settings with temporal
dependence, nonlinear structure, and heterogeneous information sources. The
four projects distinguish the quantity being estimated from the uncertainty
statement and empirical evidence used to support it.

The first project develops source-aware correction and synthesis of river-flow
forecasts. A shared latent quantile process is combined with source-specific
discrepancies, transfer covariates, and predictive synthesis. In a retrospective
five-origin San Lorenzo River study, the selected multivariate
extended-asymmetric-Laplace model gives the lowest 28-day continuous ranked
probability score among the evaluated Bayesian variants at every origin. A
separate eight-day comparison shows that relative performance varies by origin,
horizon, and comparator.

The second project develops the \pkg{exdqlm} \proglang{R} package for dynamic and
static Bayesian quantile analysis. It connects model specification, Markov chain
Monte Carlo and variational inference, diagnostics, forecasting, predictive
synthesis, and shrinkage regression in a common fitted-object interface. The
examples establish the scope of the software and show that computational speed,
posterior uncertainty, and held-out accuracy need not rank methods in the same
order.

The third project combines fixed deep echo-state-network features with Bayesian
single-level and multiquantile regression. The multiquantile formulation reduces
crossing but does not improve forecast scores uniformly. In the retrospective
applications, Q--DESN under the asymmetric Laplace model gives the lowest grid
score at one GloFAS origin, and region-specific model selection lowers mean
average quantile loss by 5.0\% relative to reproduced PriceFM forecasts over 114
region--fold cases. These comparisons are conditional on the recorded designs
and information sets.

The fourth project characterizes a mean-preserving fixed-content interval and
embeds it in a family indexed by retained-mean tilt. This geometry motivates a
profiled tolerance procedure that searches over fitted content and placement and
uses a direct posterior content check. Independent simulations provide
repeated-sampling evidence on eight distributions, while a pharmaceutical batch
analysis illustrates endpoint placement and data limitations. Regression and
dynamic endpoint formulations are stated as proposed extensions.

Together, the projects show that similar conditional-Gaussian calculations may
support different inferential claims. Likelihood-based quantile inference,
predictive synthesis, loss-based generalized updating, and tolerance calibration
must therefore be interpreted according to their respective targets and
evaluation designs.
\end{abstract}
